%
% Spezielle komplexe Funktionen
%
\subsubsection{Spezielle komplexe Funktionen}

%
% produkt: Produktdarstellungen
%
Polynome sind bis auf einen Vorfaktor durch die Nullstellen
festgelegt.
Ob dies auch für beliebige ganze Funktionen gilt, diskutieren
\emph{Simon Köpfli}
\index{Kopfli, Simon@Köpfli, Simon}%
\index{Simon Kopfli@Simon Köpfli}%
und
\emph{Florian von Wyl}
\index{von Wyl, Florian}%
\index{Florian von Wyl}%
in Kapitel~\ref{chapter:produkt}.
Da es ganze Funktionen mit unendlich vielen Nullstellen gibt, muss
man dazu unendlichen Produkten einen Sinn geben.
\index{Produkt, unendlich}%
Dabei treten zunächst Konvergenzfragen als Hindernis auf, die durch
die weierstraßschen Elementarfaktoren beantwortet werden.
Dies genügt aber nicht.
Die Exponentialfunktion $e^{g(z)}$ einer ganzen Funktion $g(z)$ hat
keine Nullstellen, die möglicherweise verschiedenen Funktionen $f(z)$
und $e^{g(z)}f(z)$ haben daher die gleichen Nullstellen, es können
aber im Allgemeinen nicht beide Funktionen Polynome sein.

%
% basel: Basel-Problem
%
Das Basel-Problem hat die besten Mathematiker des 17.~und 18.~Jahrhunderts
herausgefordert.
\index{Basel-Problem}%
Erst Leonhard Euler fand die korrekte Summe der reziproken Quadrate der
\index{Euler, Leonhard}%
natürlichen Zahlen.
\emph{Jan Gachnang}
und
\emph{Etienne Schafflützel}
präsentieren in Kapitel~\ref{chapter:basel} Eulers klassische Lösung,
die allerdings modernen Standards nicht genügt.
Strenge Beweise werden ebenfalls dargelegt.

%
% gamma: Gamma-Funktion
%
Die Fakultät ist die prototypische rekursiv definierte Funktion, sie
tritt in der Taylor-Reihe und in der Kombinatorik immer wieder auf,
und sie wird auch in den Kapiteln~\ref{chapter:zeta} und \ref{chapter:hankel}
verwendet.
\emph{Lukas Buchli}
\index{Lukas Buchli}%
\index{Buchli, Lukas}%
und
\emph{Roman Weber}
\index{Roman Weber}%
\index{Weber, Roman}%
zeigen, wie die gleiche rekursive Struktur auch in den Volumina 
geradedimensionaler Kugeln auftritt.
Sie leiten die Formel
\[
V_n(r)
=
\frac{\pi^{n/2}}{(n/2)!}r^n
\]
für das Volumen der $n$-dimensionalen Kugel her.
Die Formel funktioniert natürlich nur für gerade $n$, doch wenn man
die Fakultät auch für beliebige reelle Zwischenwerte definieren könnte,
würde die Formel auch für ungeraddimensionale Kugeln funktionieren.
Wie der Satz von Bohr-Mollerup diese Erweiterung zur Gamma-Funktion
\index{Gamma-Funktion}%
möglich macht, zeigen sie in Kapitel~\ref{chapter:gamma}.

%
% bessel: Bessel-Funktionen
%
Die Bessel-Funktionen sind Lösungen einer komplexen linearen
Differentialgleichung mit Koeffizienten, die am Rand des
Definitionsbereichs singulär sind.
Ihr vollständiges Verständnis ist nur als komplexe Funktionen
\emph{Gian Cavegn}s
\index{Gian Cavegn}%
\index{Cavegn, Gian}%
Kapitel~\ref{chapter:bessel} gibt eine erste Einführung in die
Lösung der besselschen Differentialgleichung.

%
% hankel: Hankel
%
Die Zeta-Funktion spielt eine zentrale Rolle in der analytischen
Zahlentheorie.
Sie wird wahlweise durch ein unendliches Produkt oder eine unendliche
Summe dargestellt.
Beides ist aber nur für Argumente in einem einen Teil der komplexen
Ebene möglich.
Ein vollständiges Bild der Funktion und vor allem die Funktion
entlang der Geraden $\frac12+it$, $t\in\mathbb{R}$, ist nur mit einer
analytischen Erweiterung möglich.
Wie diese im Prinzip konstruiert werden kann, wurde in
Kapitel~\ref{chapter:analytisch} diskutiert.
Die dort vorgeschlagene Methode ist aber nicht immer praktikabel.
\emph{Jannis Gull}
\index{Jannis Gull}%
\index{Gull, Jannis}%
stellt in Kapitel~\ref{chapter:hankel} eine alternative Möglichkeit
mit komplexer Integration vor, die einen leichteren Zugang zur
analytischen Erweiterung von $\zeta(s)$ ermöglicht.

%
% zeta: Die Summe der natürlichen Zahlen
%
Die Berechnung von Werten der Zeta-Funktion ist nicht immer einfach.
Euler berechnet für das Basel-Problem erfolgreich den Wert
$\zeta(2)=\frac{\pi^2}{6}$.
Kapitel~\ref{chapter:zeta} von
\emph{Melissa Haumüller}
\index{Melissa Haumüller}%
\index{Haumüller, Melissa}%
zeigt, wie sich mit der analytischen Erweiterung der $\zeta$-Funktion
der auf den ersten Blick unsinnige Wert
$\zeta(-1) = 1 + 2 + 3 +\dots = -\frac{1}{12}$ recthfertigen lässt.
